% TOP_PROBLEM: {"id": 30006877, "title": "Arithmeticity of Lattices in PU(n,1) for n at Least 4", "category_id": 17, "category": {"id": 17, "name": "group_theory", "display_name": "Group Theory", "description": "Problems about groups, group actions, representations, and related algebraic structures.", "slug": "group-theory", "order_index": 17, "created_at": "2026-07-31T15:26:25.671Z"}, "status": "open"}
% FIRST_POSTED: 2026-09-27
% !TeX program = lualatex
% ATTEMPT_STATUS: unresolved
\documentclass[11pt]{article}
\usepackage[margin=25mm]{geometry}
\usepackage{fontspec}
\setmainfont{Latin Modern Roman}
\usepackage{amsmath,amssymb,mathtools}
\usepackage{unicode-math}
\setmathfont{Latin Modern Math}
\usepackage{xurl,hyperref}
\hypersetup{colorlinks=true,urlcolor=blue}
\setcounter{secnumdepth}{0}
\setlength{\parindent}{0pt}
\setlength{\parskip}{5pt}
\setlength{\emergencystretch}{3em}
\begin{document}
\section*{\texorpdfstring{Arithmeticity of Lattices in PU(n,1) for n at Least 4}{Arithmeticity of Lattices in PU(n,1) for n at Least 4}}\label{30006877}
\subsection{Definitions and mathematical statement}
Complex hyperbolic $n$-space is the space of negative complex lines for a Hermitian form of signature $(n,1)$; its holomorphic isometry group is $\operatorname{PU}(n,1)$. A lattice $\Gamma$ is discrete and has finite Haar covolume. Two subgroups are commensurable if their intersection has finite index in both. In the arithmeticity convention here, $\Gamma$ is commensurable, after conjugation, with the projection of an arithmetic subgroup of a ${\mathbb{Q}}$-algebraic group through a real-group homomorphism with compact kernel; finite central lifts to $\operatorname{SU}(n,1)$ give the equivalent formulation. The assertion is
\[
 n\ge4,\quad\Gamma<\operatorname{PU}(n,1)\text{ a lattice}
 \quad\Longrightarrow\quad\Gamma\text{ arithmetic}.
\]
Here $n$ is complex dimension. Noncompact finite-volume quotients, cocompact lattices, and torsion are all allowed, without restricting the arithmetic construction to a particular type of Hermitian form.
\par\smallskip\textit{Formulation and concept references:} \href{https://www.maths.dur.ac.uk/users/j.r.parker/img/Lattices.pdf}{[S1]}.

\subsection{Short English statement}
Is every finite-covolume lattice in the holomorphic isometry group of complex hyperbolic space arithmetic once the complex dimension is at least four?
\subsection{Research attempt}
\textbf{Status: unresolved.} The dimension in this assertion is complex dimension. Increasing it to four does not make the ambient group a higher-real-rank group. The source [S1] asks whether nonarithmetic examples exist in these dimensions; its discussion does not establish the opposite assertion in the selected formulation. I retain that formulation as the target, without assigning the survey's author a conjecture he does not make.

The source audit included the primary paper of Bader--Fisher--Miller--Stover [E1], whose final arXiv version is dated 22 February 2023. Its Theorem 1.1 proves arithmeticity under an additional hypothesis of infinitely many maximal properly immersed totally geodesic submanifolds of real dimension at least two. That hypothesis is absent here. Its Theorem 1.5 provides an arithmetic description and integrality for complex hyperbolic lattices and reduces arithmeticity to compactness at the other archimedean places. Neither statement by itself supplies the conclusion for all complex dimensions $n\ge4$. The primary sources inspected on 27 September 2026 did not provide a resolution of that unrestricted assertion.

The work below proves several elementary reductions and constructs arithmetic examples in every relevant dimension. The standard arithmetic lattice, compactness, and density theorems are identified when used. The central gap is a universal exclusion of additional noncompact conjugate factors, not construction of another arithmetic example.

\subsubsection{1. The ambient real rank is one in every dimension}
Work first in $G=\operatorname{SU}(n,1)$ for the form
\[
 J=\operatorname{diag}(1,\ldots,1,-1).
\]
On a plane consisting of one positive and the negative coordinate, the matrices
\[
 a_t=\begin{pmatrix}\cosh t&\sinh t\\
                         \sinh t&\cosh t\end{pmatrix}
 \tag{351.1}
\]
preserve $\operatorname{diag}(1,-1)$ and have determinant one. Extending by the identity gives a one-dimensional real split torus: in an isotropic basis its two nontrivial eigenvalues are $e^t,e^{-t}$. Thus the real rank is at least one.

For completeness, the upper bound follows directly from the Hermitian form. Let $T$ be a real split algebraic torus of dimension $r$ in $G$. On the complex vector space, its representation decomposes into character spaces $V_\chi$. On the identity component of its real points, the characters have positive real values. If $v\in V_\chi$ and $w\in V_\psi$, invariance gives
\[
 h(v,w)=\chi(t)\psi(t)h(v,w)\quad(t\in T(\mathbb R)^0).
\]
Therefore $h(V_\chi,V_\psi)=0$ unless $\chi\psi=1$. In additive notation, nonzero weights are paired with their negatives. Choose a generic linear functional on the real span of the character lattice, nonzero on every nonzero weight. The sum of the weight spaces on its positive side is totally isotropic, since the sum of two positive weights cannot be zero.

A Hermitian space of signature $(n,1)$ has totally isotropic subspaces of dimension at most one. Indeed, projection from such a subspace onto the negative coordinate is injective: a vector in its kernel lies in a positive-definite subspace and cannot be isotropic unless zero. It follows that there is at most one positive nonzero weight, counted with multiplicity. But the weights span a space of rank $r$, because the representation of $T$ is faithful. Their pairing then forces $r\le1$. Together with (351.1), this proves
\[
 \operatorname{rank}_{\mathbb R}\operatorname{SU}(n,1)=1.
 \tag{351.2}
\]
The finite central quotient to $\operatorname{PU}(n,1)$ does not change this rank.

Thus a proof beginning ``$n\ge4$, so apply higher-rank arithmeticity'' fails at its first hypothesis. Matrix size, complex dimension, and real rank are different parameters. The existence of large compact unitary subgroups likewise does not increase the real split rank.

\subsubsection{2. Finite central lifts and commensurability do not change the target}
Let $q:\widetilde G\to G$ be a surjective Lie-group homomorphism with finite central kernel. If $\Gamma<G$ is discrete, then $q^{-1}(\Gamma)$ is discrete: use a neighborhood on which $q$ is a local homeomorphism and which meets no nonidentity kernel element. Conversely, the image of a discrete subgroup of $\widetilde G$ is discrete. To see the converse carefully, a convergent sequence of distinct image elements has lifts approaching one of finitely many points in a compact inverse image of a small compact neighborhood; this contradicts discreteness. Finite coverings compare Haar measures by a fixed finite factor. Hence finite covolume is preserved in both directions.

For $\operatorname{SU}(n,1)\to\operatorname{PU}(n,1)$ the kernel consists of the scalar matrices of determinant one and has order $n+1$. The arithmeticity convention in the root statement was chosen to be compatible with this lift. We may therefore work upstairs without assuming the lattice is torsion free or losing the noncompact finite-volume case.

Commensurability is transitive. If $\Gamma_1\cap\Gamma_2$ has finite index in both groups, and so does $\Gamma_2\cap\Gamma_3$, then their intersection has finite index in $\Gamma_2$. Intersecting that subgroup with each $\Gamma_i$ shows that $\Gamma_1\cap\Gamma_3$ has finite index in both $\Gamma_1$ and $\Gamma_3$. More explicitly, intersections of two finite-index subgroups have index at most the product of the two indices; this proves the needed estimates. Consequently replacing a lattice by a finite-index subgroup preserves arithmeticity in both directions.

These reductions are useful only when the relevant finite-index or finite-kernel fact has been proved. Passing to an arbitrary infinite-index subgroup of an arithmetic lattice preserves containment but need not preserve finite covolume. The next observation isolates that issue.

\subsubsection{3. A lattice contained in a discrete arithmetic group has finite index}
\textbf{Lemma 351.1.} Suppose $\Gamma\le\Lambda\le G$, both are discrete, and $\Gamma$ is a lattice. Then $[\Lambda:\Gamma]<\infty$. In particular, if $\Lambda$ is arithmetic, then $\Gamma$ is arithmetic.

\emph{Proof.} Choose a relatively compact identity neighborhood $U$ of positive Haar measure so small that distinct left translates $\lambda U$ with $\lambda\in\Lambda$ are disjoint. Such a neighborhood exists by discreteness, after requiring $UU^{-1}\cap\Lambda=\{1\}$. For each left coset $\Gamma\lambda$ in $\Lambda$, its image $\Gamma\lambda U$ in $\Gamma\backslash G$ has measure $\mu(U)$. These images for different cosets are disjoint. Indeed, an equality $\gamma\lambda u=\lambda' u'$ would give a nontrivial element of $\Lambda$ in $UU^{-1}$ unless $u=u'$ and $\Gamma\lambda=\Gamma\lambda'$. Infinitely many cosets would therefore force infinite quotient measure, contradicting the lattice assumption. The final assertion follows from the definition of commensurability. $\square$

When both covolumes are already known, the same decomposition gives
\[
 \operatorname{vol}(\Gamma\backslash G)
  =[\Lambda:\Gamma]\operatorname{vol}(\Lambda\backslash G).
 \tag{351.3}
\]
Thus a genuinely sufficient certificate for arithmeticity is containment, after conjugation and finite-index passage, in a \emph{discrete} arithmetic lattice in the same group. Merely producing integral matrices in a larger product of noncompact groups does not give this certificate: the projection of that arithmetic group to one factor can be nondiscrete.

\subsubsection{4. Explicit arithmetic examples in every complex dimension}
First take $E=\mathbb Q(i)$ and
\[
 h_0(z,z)=|z_1|^2+\cdots+|z_n|^2-|z_{n+1}|^2.
\]
The determinant-one $\mathbb Z[i]$-linear automorphisms preserving this form give an arithmetic subgroup $\Lambda_0$ of $\operatorname{SU}(n,1)$. Discreteness follows directly from the discreteness of Gaussian integral matrix entries. To conclude finite covolume, one uses the standard Borel--Harish-Chandra theorem: integral points in a connected semisimple group defined over $\mathbb Q$ form an arithmetic lattice in its real points. This is a cited general input, not an elementary consequence of discreteness. The unitary group here is defined over $\mathbb Q$ through the quadratic extension $E/\mathbb Q$, so the theorem applies. Projection gives an arithmetic lattice in the group of the problem.

There is also an instructive construction with a compact extra real factor. Let
\[
 K=\mathbb Q(\sqrt2),\qquad E=K(i),\qquad
 h_1=\operatorname{diag}(1,\ldots,1,-\sqrt2).
 \tag{351.4}
\]
At the embedding $\sigma_+:\sqrt2\mapsto+\sqrt2$, its signature is $(n,1)$. At $\sigma_-:\sqrt2\mapsto-\sqrt2$, it is positive definite. Write $\mathcal O_E$ for the full ring of integers and take the determinant-one form-preserving automorphisms of $\mathcal O_E^{n+1}$. Under the two archimedean embeddings they lie diagonally in
\[
 \operatorname{SU}(n,1)\times\operatorname{SU}(n+1),
 \tag{351.5}
\]
up to a real change of basis in each factor. Restriction of scalars supplies a semisimple $\mathbb Q$-group with these real points. The same arithmetic lattice theorem gives a lattice $\widetilde\Lambda_1$ in the product.

Why is its projection to the first factor discrete? Suppose projected elements converge to the identity. The second coordinates lie in a compact group, so a subsequence converges. Quotients of successive distinct elements of that subsequence would converge to the identity in the whole product, violating discreteness of $\widetilde\Lambda_1$. Thus the identity is isolated in the projected group. Finite covolume also descends through the compact factor: integration over that factor, with Haar measure normalized to one, identifies the quotient measures up to a fixed normalization. Equivalently, the natural map from the finite-volume product quotient to the first-factor quotient has compact homogeneous fibers and pushes its invariant measure to a finite nonzero invariant measure.

This construction is cocompact by the standard anisotropy compactness criterion for arithmetic groups. The relevant anisotropy can be checked here: if a nonzero vector $v\in E^{n+1}$ satisfied $h_1(v,v)=0$, applying $\sigma_-$ would give a zero value of a positive-definite Hermitian form on a nonzero vector, which is impossible. For special unitary groups this absence of an isotropic vector is the absence of a nontrivial $K$-split torus. The compactness criterion then applies to the restriction-of-scalars group. The proof of that criterion is not supplied by the signature calculation; the calculation verifies its hypothesis.

Both families exist for every $n\ge4$ and fit the allowed compact-kernel arithmeticity convention. They demonstrate consistency of the desired conclusion with many examples. They do not classify an arbitrary lattice, and the use of Hermitian forms in these examples does not narrow the universal target to arithmetic groups of this construction type.

\subsubsection{5. The conjugate-place obstruction can be isolated precisely}
Here is the specific external structural input that sharpens the remaining question. In the setting of [E1, Theorem 1.5], a lattice in $\operatorname{SU}(n,1)$ has totally real adjoint trace field $K$ and an associated simply connected $K$-group $\mathbf G$. At each real place $v$ its real points have type $\operatorname{SU}(r_v,s_v)$ with $r_v+s_v=n+1$. Integrality holds after finite-index passage in an appropriate representation. Arithmeticity is equivalent to all archimedean factors other than the distinguished lattice factor being compact. These are imported structural conclusions, not proved by the examples above.

A direct consequence, conditional only on that established structural theorem, is the following useful subcase:
\[
 K=\mathbb Q\quad\Longrightarrow\quad\Gamma\text{ is arithmetic}.
 \tag{351.6}
\]
There is only one real place, so the compactness condition at the other archimedean places is vacuous. More generally, the same argument works whenever each other signature is definite. This identifies a concrete route for proving arithmeticity of an individual lattice: determine its associated field and group, and verify the remaining signatures.

For the universal assertion, however, knowing $n\ge4$ does not force $s_v=0$ or $r_v=0$ at those other places. The elementary inequality $r_v+s_v=n+1$ allows many indefinite signatures. A proof must use the fact that the distinguished representation has discrete finite-covolume image to eliminate those possibilities, or find an independent arithmeticity mechanism. There is no such elimination in the rank computation or the form examples.

This discussion also refines the slogan that arithmeticity is ``integrality plus a number field.'' For the actual complex hyperbolic lattice, the structural theorem already supplies integrality in an appropriate model. The remaining obstruction is an archimedean one. For a collection of arbitrary matrices, even integrality and discreteness first have to be proved; for a genuine lattice, repeatedly proving those weaker properties cannot substitute for compactness of the other factors.

\subsubsection{6. Two invalid ways to manufacture a counterexample}
\emph{Embedding a smaller-dimensional lattice.} Put a lattice in $\operatorname{SU}(m,1)$ into $\operatorname{SU}(n,1)$ by adjoining identity blocks, with $m<n$. Its matrices fix a nonzero positive subspace pointwise. They lie in a proper real algebraic subgroup of the larger special unitary group. The Borel density theorem says that a lattice in this connected semisimple group with no compact factors is Zariski dense. Thus the block-embedded group cannot be a lattice of the larger group. Discreteness survives this embedding, but finite covolume does not. In particular, a nonarithmetic example in complex dimension two or three cannot simply be enlarged this way to answer the selected question negatively.

The same dimension check applies to a real hyperbolic lattice preserving a real quadratic form of signature $(n,1)$. Regarding its real matrices as complex matrices places it in a proper real algebraic subgroup of $\operatorname{SU}(n,1)$. A lattice for real hyperbolic $n$-space is not thereby a lattice for complex hyperbolic $n$-space. A real dimension-four construction is especially easy to misread when the selected parameter is complex dimension four.

\emph{Algebraic entries alone.} Let
\[
 \alpha=(3+4i)/5,
\]
and take a diagonal matrix $g$ with entries $\alpha$ and $\alpha^{-1}$ on one positive and the negative coordinate and entry one elsewhere. Then $g^*Jg=J$, $\det g=1$, and all entries lie in $\mathbb Q(i)$. Nevertheless $\langle g\rangle$ is nondiscrete.

To verify the last claim, first observe that $\alpha$ is not a root of unity. If it were, both $\alpha$ and $\overline\alpha$ would be algebraic integers. Their sum $6/5$, being rational, would be an integer, a contradiction. Partition the unit circle into $N$ equal arcs. Among $1,\alpha,\ldots,\alpha^N$ two lie in one arc, so a nonzero power $\alpha^{q_N}$ is within $2\pi/N$ of one. These powers are not one, and a sequence of distinct such powers approaches one; otherwise a finite nonempty set of distances from one would have positive minimum. The corresponding distinct powers of $g$ approach the identity. This is a fully explicit failure of the implication ``matrices over a number field imply a lattice,'' before arithmeticity is even considered.

\subsubsection{7. Checks, scope, and the first missing argument}
The rational parametrization
\[
 a=\frac{t^2+1}{t^2-1},\qquad b=\frac{2t}{t^2-1},\qquad
 a^2-b^2=1\quad(t\ne\pm1)
 \tag{351.7}
\]
provides exact matrix checks for the two-coordinate form-preserving calculation. The diagnostic accompanying this attempt verifies these identities over rational numbers and checks the Gaussian rational norm and trace used for $\alpha$. Such checks can catch a sign or determinant error. They cannot certify finite covolume; in the constructions that property came from the arithmetic lattice theorem with its hypotheses explicitly verified.

The proved elementary reductions are the rank-one calculation, finite-index and finite-central invariances, the containment criterion, the signature checks for the example families, and the nondiscrete algebraic matrix example. The arithmetic construction theorems, density theorem, and structural theorem [E1] are clearly separate imported inputs. Their consequences include arithmeticity when the associated trace field is $\mathbb Q$, and when all conjugate archimedean forms are compact.

The first missing universal step is to show that every lattice of complex dimension at least four satisfies that compactness condition, or to produce a verified finite-covolume lattice violating it. The extra geodesic hypothesis of [E1, Theorem 1.1] cannot simply be deleted; the construction of arithmetic examples cannot be promoted to an exhaustive classification; and higher-rank arithmeticity cannot be invoked in a real-rank-one group. The full assertion in the root statement remains unresolved by this attempt.

\subsection{Sources}
\begin{itemize}
\item[S1]Complex hyperbolic lattices; undated author-hosted survey snapshot, not assigned a publication date from PDF metadata.
\url{https://www.maths.dur.ac.uk/users/j.r.parker/img/Lattices.pdf}.
\textit{Formulation source.}
\item[E1]Uri Bader, David Fisher, Nicholas Miller, and Matthew Stover, \emph{Arithmeticity, superrigidity and totally geodesic submanifolds of complex hyperbolic manifolds}, arXiv:2006.03008v2, 22 February 2023, published in \emph{Inventiones Mathematicae} 233 (2023), 169--222.
\url{https://arxiv.org/abs/2006.03008}.
\textit{Added primary source. Theorems 1.1 and 1.5 and the introduction inspected 27 September 2026; the extra geodesic hypothesis and compactness-at-other-places criterion are retained.}
\end{itemize}
\end{document}
